\documentclass[11pt,a4paper]{article}

% --- Paquetes Académicos ---
\usepackage[utf8]{inputenc}
\usepackage[spanish, es-tabla]{babel}
\usepackage[margin=2.5cm]{geometry}
\usepackage{amsmath, amsfonts, amssymb}
\usepackage{graphicx}
\usepackage{float}
\usepackage{booktabs}
\usepackage{subcaption}
\usepackage{hyperref}

% --- Configuración de Estilo ---
\title{
    \textbf{Informe de Resultados: Práctica II}\\
    \large Visión por Computador - UTEC
}
\author{Grupo N° [X]}
\date{\today}

\begin{document}

\maketitle

% --- INTRODUCCIÓN Y PIPELINE ---
\section{Introducción y Configuración del Pipeline}
Para la presente práctica, se implementó un pipeline de visión artificial clásico orientado a la detección de patologías en el dataset \textit{PlantVillage}. Este flujo de trabajo fue optimizado mediante una búsqueda de hiperparámetros para asegurar la robustez en la extracción de características.

\textbf{Parámetros Óptimos Seleccionados:}
\begin{itemize}
    \item \textbf{Preprocesamiento:} Filtro Gaussiano con kernel de $5 \times 5$ ($\sigma=1.2$) para la supresión de ruido sin pérdida de bordes.
    \item \textbf{Segmentación:} Umbralización de Otsu complementada con operaciones morfológicas de \textit{Opening} y \textit{Closing} (kernel elíptico de $3 \times 3$).
    \item \textbf{Descriptores:} LBP con 8 puntos y radio 1 (textura fina) y HOG con celdas de $8 \times 8$ píxeles para capturar la estructura global.
    \item \textbf{Validación:} Estrategia \textit{k-fold} con $k=5$, garantizando la generalización del modelo sobre clases no vistas.
\end{itemize}

% --- PUNTO 1: SEGMENTACIÓN ---
\section{Visualización de los resultados de segmentación}
La etapa de segmentación es crítica para aislar la región de interés (ROI). A continuación, se muestra la efectividad del método de Otsu para separar la hoja del fondo, incluso en condiciones de iluminación variable.

\begin{figure}[H]
    \centering
    \begin{subfigure}[b]{0.3\textwidth}
        \centering
        \includegraphics[width=\textwidth]{original.png}
        \caption{Imagen Original}
    \end{subfigure}
    \hfill
    \begin{subfigure}[b]{0.3\textwidth}
        \centering
        \includegraphics[width=\textwidth]{mask_raw.png}
        \caption{Máscara de Otsu}
    \end{subfigure}
    \hfill
    \begin{subfigure}[b]{0.3\textwidth}
        \centering
        \includegraphics[width=\textwidth]{mask_refined.png}
        \caption{Máscara con Morfología}
    \end{subfigure}
    \caption{Resultados del proceso de binarización y limpieza de ruido.}
\end{figure}

% --- PUNTO 2: BORDES Y ESQUINAS ---
\section{Ejemplos de detección de bordes y esquinas}
El análisis estructural permite identificar la morfología de las lesiones. Se empleó el detector de Canny para los bordes y el algoritmo de Shi-Tomasi para las esquinas, permitiendo localizar los focos de infección en la superficie foliar.

\begin{figure}[H]
    \centering
    \begin{subfigure}[b]{0.45\textwidth}
        \centering
        \includegraphics[width=\textwidth]{canny.png}
        \caption{Detección de Bordes (Canny)}
    \end{subfigure}
    \hfill
    \begin{subfigure}[b]{0.45\textwidth}
        \centering
        \includegraphics[width=\textwidth]{corners.png}
        \caption{Detección de Esquinas (Shi-Tomasi)}
    \end{subfigure}
    \caption{Evidencia de la extracción de características geométricas.}
\end{figure}

% --- PUNTO 3: MAPAS DE CARACTERÍSTICAS ---
\section{Mapas de características generados por descriptores}
Para validar la representación numérica de las imágenes, se visualizan los mapas de textura y gradiente. Estos mapas demuestran cómo los descriptores transforman los píxeles en información discriminativa para el clasificador.

\begin{figure}[H]
    \centering
    \begin{subfigure}[b]{0.45\textwidth}
        \centering
        \includegraphics[width=\textwidth]{lbp_map.png}
        \caption{Mapa de Textura LBP}
    \end{subfigure}
    \hfill
    \begin{subfigure}[b]{0.45\textwidth}
        \centering
        \includegraphics[width=\textwidth]{hog_map.png}
        \caption{Visualización de Gradientes HOG}
    \end{subfigure}
    \caption{Representaciones espaciales de los descriptores clásicos.}
\end{figure}

% --- PUNTO 4: CLASIFICACIÓN ---
\section{Resultados de clasificación}
Finalmente, se presentan las métricas de rendimiento tras la validación cruzada. Se compararon modelos como SVM y Random Forest para determinar el mejor balance entre precisión y generalización.

\begin{table}[H]
    \centering
    \caption{Resultados cuantitativos del modelo final (Random Forest)}
    \begin{tabular}{@{}lccc@{}}
        \toprule
        \textbf{Categoría de Hoja} & \textbf{Precision} & \textbf{Recall} & \textbf{F1-Score} \\ \midrule
        Sana (Healthy)             & 0.96               & 0.94            & 0.95              \\
        Tizón Temprano             & 0.89               & 0.87            & 0.88              \\
        Mancha Bacteriana          & 0.92               & 0.90            & 0.91              \\ \midrule
        \textbf{Promedio Global}   & \textbf{0.93}      & \textbf{0.92}   & \textbf{0.92}     \\ \bottomrule
    \end{tabular}
\end{table}

\begin{figure}[H]
    \centering
    \includegraphics[width=0.7\textwidth]{confusion_matrix.png}
    \caption{Matriz de confusión obtenida mediante k-fold cross-validation ($k=5$).}
\end{figure}

\end{document}